% !PRINTED 319
\weberpage{319}
%<<CONT
8\textsuperscript{ten} Grades sind, worauf hier nicht näher eingegangen
werden soll (vgl. Bd. I, \S.\,185).

\begin{center}
\S.\ 84.

\gesperrt{Biquadratische Gleichungen.}
\end{center}

Wir haben uns auf die Betrachtung der primitiven metacyklischen
Gleichungen 8\textsuperscript{ten} Grades beschränkt, weil die imprimitiven
auf Gleichungen 4\textsuperscript{ten} Grades zurückkommen. Um also auch die
Wurzeln der letzteren in der Weise des vorigen Paragraphen darstellen zu
können, haben wir noch eine analoge Darstellung der Wurzeln biquadratischer
Gleichungen zu suchen. Hierbei sind dieselben Betrachtungen, wie im
\S.\,82, nur in wesentlich einfacherer Gestalt anwendbar.

Wir bezeichnen die Wurzeln der biquadratischen Gleichung mit
\weq{1}{X_{0,0}, X_{1,0}, X_{0,1}, X_{1,1},}

und können dann die ganze Permutationsgruppe 24\textsuperscript{sten}
Grades als binäre lineare Congruenzgruppe für den Modul 2 darstellen:
\weq{2}{P = SQ.}

Die darin enthaltene homogene Gruppe $S$ ist vom 6\textsuperscript{ten}
Grade und besteht aus den Substitutionen
\weq{3}{\sigma = \begin{pmatrix}1, & 0\\0, & 1\end{pmatrix},
\begin{pmatrix}1, & 0\\1, & 1\end{pmatrix},
\begin{pmatrix}0, & 1\\1, & 0\end{pmatrix},
\begin{pmatrix}0, & 1\\1, & 1\end{pmatrix},
\begin{pmatrix}1, & 1\\1, & 0\end{pmatrix},
\begin{pmatrix}1, & 1\\0, & 1\end{pmatrix},}

und man kann
\weq{4}{\sigma_1 = \begin{pmatrix}0, & 1\\1, & 0\end{pmatrix}, \quad
\sigma_2 = \begin{pmatrix}1, & 1\\0, & 1\end{pmatrix}}

als die erzeugenden Substitutionen der Gruppe $S$ ansehen.

Wir definiren nun wie im \S.\,80 die drei Grössen:
\weq{5}{\begin{aligned}
\Psi_{1,0} &= X_{0,0} - X_{1,0} + X_{0,1} - X_{1,1}\\
\Psi_{0,1} &= X_{0,0} + X_{1,0} - X_{0,1} - X_{1,1}\\
\Psi_{1,1} &= X_{0,0} - X_{1,0} - X_{0.1} + X_{1,1}, % [sic] printed with a period after 0 instead of a comma
\end{aligned}}

so dass $\Psi^2_{1,0}, \Psi^2_{0,1}, \Psi^2_{1,1}$, aber auch das Product
$\Psi_{1,0}\Psi_{0,1}\Psi_{1,1}$ durch die Substitutionen der cyklischen
Gruppe $Q$ ungeändert bleiben.
